\chapter{Estensione agentica}
\label{chap:estensione-agentica}

Il presente capitolo copre il secondo blocco di attività dello stage (settimane 5--8 del Piano di lavoro), dedicato allo studio e all'eventuale integrazione di sistemi agentici nell'Accessibility Checker sviluppato nei capitoli precedenti. Diversamente dal blocco principale (capitoli~\ref{chap:metodologia-selezione}--\ref{chap:validazione-benchmark}), questo blocco ha carattere \emph{desiderabile} ed è subordinato alla buona riuscita del blocco principale, come indicato nel Piano di lavoro approvato dalla Prof.ssa Gaggi: le sezioni seguenti documentano quanto effettivamente realizzato in questa fase.

\section{Studio di Pi Agent e OpenCode}
\label{sec:studio-pi-agent-opencode}

\todoplaceholder{studio dei sistemi Pi Agent e OpenCode e analisi delle possibilità di utilizzo di sistemi agentici nel progetto.}

% TODO (settimana 5 del piano di lavoro): studio di Pi Agent e di OpenCode
% come possibili framework/riferimenti per un'estensione agentica del
% sistema; analisi delle possibilità concrete di utilizzo di sistemi agentici
% nel contesto dell'Accessibility Checker, alla luce dei criteri \emph{core} e del
% prototipo già sviluppati.


\section{Progettazione dell'estensione agentica}
\label{sec:progettazione-estensione-agentica}

\todoplaceholder{progettazione dell'estensione agentica del sistema.}

% TODO (settimana 5 del piano di lavoro): progettazione dell'estensione
% agentica, comprensiva delle agentic skill previste, dell'integrazione con
% il sistema sviluppato nel blocco principale e delle scelte architetturali
% conseguenti.


\section{Implementazione delle agentic skill}
\label{sec:implementazione-agentic-skill}

\todoplaceholder{implementazione delle agentic skill e integrazione con il sistema sviluppato nel primo blocco.}

% TODO (settimana 6 del piano di lavoro): implementazione di almeno una
% agentic skill (obiettivo desiderabile) ed eventuali skill aggiuntive
% (obiettivo opzionale); integrazione con il sistema del blocco principale;
% test.


\section{Valutazione del sistema completo e confronto con l'approccio basato su prompt engineering}
\label{sec:valutazione-approccio-agentico}

\todoplaceholder{valutazione del sistema completo, confronto tra approccio basato su prompt engineering e approccio agentico, valutazione del grado di automazione ottenibile.}

% TODO (settimana 7 del piano di lavoro): valutazione del sistema completo,
% comprensivo dell'estensione agentica; confronto sistematico tra l'approccio
% basato su prompt engineering (capitoli~\ref{chap:progettazione-prototipo}
% e~\ref{chap:validazione-benchmark}) e l'approccio agentico qui descritto;
% valutazione del grado di automazione effettivamente ottenibile con
% quest'ultimo.


\section{Consolidamento, test finali e sviluppi futuri}
\label{sec:consolidamento-agentico}

\todoplaceholder{consolidamento del lavoro, test finali, e possibili sviluppi futuri dell'estensione agentica.}

% TODO (settimana 8 del piano di lavoro): consolidamento del lavoro svolto
% sull'estensione agentica; test finali; possibili sviluppi futuri specifici
% di questo blocco (ad esempio un workflow agentico più esteso o
% un'ulteriore validazione, entrambi obiettivi opzionali). La documentazione
% tecnica conclusiva dell'intero progetto è presentata nel
% capitolo~\ref{chap:conclusioni}.
